\chapter{Formation initiale et diplômes}

\section{Parcours scolaire initial}

\subsection{Baccalauréat Sciences et Technologies de Laboratoire (1996)}

Mon parcours dans l'enseignement supérieur technique débute par l'obtention d'un Baccalauréat Sciences et Technologies de Laboratoire (STL) option Chimie de Laboratoire en 1996. Ce diplôme m'a apporté une formation scientifique solide, développant des compétences en analyse, rigueur méthodologique et esprit logique, compétences transversales essentielles pour l'informatique \parencite{competences_21_siecle}.

\subsection{DUT Informatique Génie Logiciel (1999)}

Poursuivant dans la voie technique, j'ai obtenu en 1999 un DUT Informatique option Génie Logiciel à l'IUT du Havre. Cette formation de deux ans m'a apporté des bases solides en :

\begin{itemize}[leftmargin=*]
    \item \textbf{Programmation} : Algorithmique, structures de données, langages (C, Java, SQL)
    \item \textbf{Génie logiciel} : Analyse, conception, méthodologies (Merise, UML)
    \item \textbf{Systèmes d'exploitation} : Architecture, administration (Unix, Windows)
    \item \textbf{Bases de données} : Modélisation relationnelle, SQL, normalisation
    \item \textbf{Réseaux} : Architecture, protocoles TCP/IP, administration
\end{itemize}

Cette formation correspond parfaitement aux attendus scientifiques du Master MEEF PIF \parencite{rncp38152}, couvrant les domaines d'ingénierie et de didactique du numérique \parencite{eduscol_nsi}.

\subsection{BTS Support Technique Micro et Réseaux (2001)}

Pour approfondir mes compétences en administration systèmes et réseaux, j'ai poursuivi avec un BTS Support Technique Micro et Réseaux obtenu en 2001. Ce diplôme complémentaire m'a permis de développer une expertise en :

\begin{itemize}[leftmargin=*]
    \item Administration de serveurs Windows et Linux
    \item Configuration et maintenance de réseaux locaux
    \item Support technique et diagnostic de pannes
    \item Sécurité informatique et gestion des sauvegardes
\end{itemize}

\section{Certifications professionnelles}

\subsection{Certifications Microsoft (2000)}

En parallèle de ma formation initiale, j'ai obtenu en autodidacte plusieurs certifications Microsoft démontrant ma capacité d'apprentissage autonome :

\begin{itemize}[leftmargin=*]
    \item \textbf{MCSE (Microsoft Certified Systems Engineer)} : Certification de haut niveau en administration de systèmes Windows Server
    \item \textbf{MCP (Microsoft Certified Professional)} : Certification de base sur les technologies Microsoft
\end{itemize}

Ces certifications valident une expertise reconnue internationalement en administration de systèmes, compétence directement transférable vers l'enseignement de l'architecture des ordinateurs et des systèmes d'exploitation en NSI.

\subsection{Certification CISCO CCNA (2000)}

J'ai également obtenu la certification \textbf{CISCO CCNA (Cisco Certified Network Associate)} en 2000, validant mes compétences en :

\begin{itemize}[leftmargin=*]
    \item Configuration de routeurs et switches CISCO
    \item Protocoles de routage (RIP, OSPF, EIGRP)
    \item Réseaux locaux virtuels (VLAN)
    \item Sécurité réseau et listes de contrôle d'accès (ACL)
\end{itemize}

Cette certification correspond aux contenus des programmes NSI terminale sur les réseaux \parencite{nsi_programmes} et démontre une maîtrise approfondie des concepts enseignés.

\section{Formation continue professionnalisante}

\subsection{Formation Concepteur Architecte en Ingénierie de Projets (2016-2018)}

En contrat de professionnalisation chez TOTAL/Stéria, j'ai suivi une formation de niveau II (Bac+4) « Concepteur Architecte en Ingénierie de Projets » validée en 2018. Cette formation de 18 mois m'a apporté :

\begin{itemize}[leftmargin=*]
    \item \textbf{Architecture logicielle} : Patterns de conception, architectures distribuées, microservices
    \item \textbf{Gestion de projet} : Méthodes agiles (Scrum, Kanban), cycle en V, DevOps
    \item \textbf{Management d'équipe} : Leadership, communication, résolution de conflits
    \item \textbf{Veille technologique} : Capacité à analyser et intégrer les innovations
\end{itemize}

Cette formation illustre ma démarche de formation continue tout au long de ma carrière, attitude essentielle pour un enseignant d'informatique face à la rapidité d'évolution de la discipline \parencite{referentiel_competences}.

\section{Diplômes d'enseignement}

\subsection{CAPES Numérique et Sciences Informatiques (2020)}

Après une reconversion maîtrisée, j'ai obtenu le CAPES NSI en 2020 \parencite{capes_nsi}, premier concours de recrutement des enseignants d'informatique en France. Ce diplôme valide :

\begin{itemize}[leftmargin=*]
    \item La maîtrise des contenus disciplinaires du programme NSI
    \item Les compétences pédagogiques et didactiques
    \item La capacité à enseigner devant un jury
    \item La connaissance du système éducatif français
\end{itemize}

\subsection{CAPET Sciences Industrielles (2020)}

J'ai également obtenu simultanément le CAPET Sciences Industrielles en 2020, démontrant une polyvalence dans l'enseignement technologique. Cette double certification constitue un atout pour enseigner en lycée technologique et professionnel.

\subsection{Master MEEF Informatique (2026)}

En juin 2026, j'ai obtenu le \textbf{Master MEEF Mention Second Degré - Parcours Informatique}. Ce diplôme universitaire de niveau Bac+5 vient consacrer ma maîtrise des fondamentaux de la didactique de l'informatique, des savoirs scientifiques disciplinaires et de la recherche en éducation numérique. Il atteste formellement de mon niveau Bac+5 académique.

\section{Synthèse de la formation initiale}

Mon parcours de formation révèle une progression cohérente vers l'expertise informatique et pédagogique :

\begin{itemize}[leftmargin=*]
    \item \textbf{Solidité technique et académique} : DUT + BTS + Master MEEF Informatique (2026) + certifications professionnelles couvrent l'ensemble des domaines du numérique.
    \item \textbf{Apprentissage autonome} : Certifications obtenues en autodidacte témoignent d'une capacité d'auto-formation.
    \item \textbf{Formation continue} : Formation Concepteur Architecte en parallèle de l'activité professionnelle.
    \item \textbf{Légitimité pédagogique et universitaire} : Master MEEF Informatique + CAPES NSI + CAPET SI valident les compétences d'enseignement et le niveau d'excellence académique.
\end{itemize}

Cette formation initiale et continue correspond parfaitement aux attendus du Master MEEF Pratiques et Ingénierie de la Formation (PIF) \parencite{rncp38152,unicaen}, tant pour les aspects disciplinaires que pour la formation par la recherche \parencite{referentiel_meef}.